To rigorously determine whether a predictive model meaningfully outperforms the random baseline (Coin Flip) or the standard reference (Engagement Index, EI), statistical significance was evaluated using the Wilcoxon signed-rank test. Unlike standard parametric tests (such as the paired t-test) which assume performance metrics are normally distributed, the Wilcoxon test is a non-parametric alternative. This makes it particularly robust for highly individual, small-sample physiological data such as EEG ($N=25$ participants).

Rather than comparing a single aggregated mean across the dataset, the test utilizes paired metrics derived from the 25 participant-specific evaluations, where an individualized model is trained and tested exclusively on each participant's own data splits. The procedure calculates the difference in performance (e.g., F1-score) between the evaluated model and the baseline for each specific participant. By ranking the absolute values of these differences and applying their original signs, the test evaluates whether the median of the differences is significantly shifted from zero. This participant-level comparison ensures that a model only achieves significance if it consistently outperforms the baseline across the majority of individuals, preventing a scenario where exceptional performance on a few participants artificially inflates the overall mean. A visual representation of these paired differences across all participants is provided in the significance heatmap (see Figure~\ref{fig:viz17_2_intra_significance}), which color-codes the exact performance gap for each participant.

The algorithm employs a two-sided Wilcoxon test with a standard significance threshold ($\alpha = 0.05$). Because a two-sided test detects significant differences in either direction (meaning the model could be significantly better \textit{or} significantly worse than the baseline), a directional validation check is applied post-hoc. A model's performance is deemed statistically significantly superior only if it meets two strict conditions: first, the test must yield $p < 0.05$; second, the arithmetic mean of the model's scores across all participants must be strictly greater than the baseline's mean. The results of this rigorous evaluation, including the final validated significance flags for each metric, are summarized in Table~\ref{tab:metrics_1_intra}. Furthermore, the test assumes that metric distributions are independent between participants, but appropriately treats models evaluated on the identical set of participant features as dependent, paired samples. Zero-differences (ties) are naturally handled by standard removal from the ranking process.
